\chapter{Bulletproofs（未完成）(Bulletproofs (WIP))}
\label{chapter:bulletproofs}
%this chapter was not finished since I decided not to include a bulletproofs explanation in ztm2


%----------
\section{向量知识证明 (Vector Knowledge Proof)}
\label{sec:vectorzkproof}

证明已知向量集合 $\boldsymbol{V}$ 中的所有元素，该集合包含 $m$ 个大小为 $N$ 的向量（$m \neq N$），其承诺为 $(C_1, C_2,..., C_m)$，其中 $C_i = v_{i,1}*H_1 + v_{i,2}*H_2 + ... + v_{i,N}*H_N$。生成元向量 $\boldsymbol{H} = \langle H_1,..., H_N \rangle$ 中各生成元关于 $G$ 的离散对数 $\boldsymbol{\lambda} = \langle \lambda_1, ..., \lambda_N \rangle$ 是未知的（例如 $H_2$ 关于 $H_4$ 的离散对数也是未知的）。验证者知道共有 $m*N$ 个元素，但无法获得关于它们的任何信息。

\subsubsection*{非交互式证明 (Non-interactive proof)}

\begin{enumerate}
	\item 生成随机整数向量（大小 N+1）\(\boldsymbol{\alpha} \in_R \mathbb{Z}_l\)，并计算内积 $C_{\alpha} = \boldsymbol{\alpha} \bullet \textbf{H}$。\footnote{向量 \(\boldsymbol{v} = \langle 1, 2, 3 \rangle\) 和 \(\boldsymbol{z} = \langle 4, 5, 6 \rangle\) 之间的内积（也称点积）为 \(\boldsymbol{v} \bullet \boldsymbol{z} = 1*4 + 2*5 + 3*6 = 32\)。}
	\item 计算{\em 挑战} $c = \mathcal{H}(T_{vec},\boldsymbol{V},C_{\alpha})$。\footnote{此处我们显式地用 $T_{vec}$ 进行域分离，用 $\boldsymbol{V}$ 进行密钥前缀。在本章其余部分这些以省略号（...）隐含。这些省略号还可以包含消息 $\mathfrak{m}$，以使证明成为签名。}
	\item 定义{\em 响应向量}（大小 N+1）$\textbf{r}$，其中对每个生成元 $H_j$ 有一个元素 $r_j$（想象向量按行排列，此处我们对每个向量 $i$ 取列求和得到一个 $v_{i,j}$）（注意每个都乘以 $c$ 的 $i$ 次幂）
	\[r_j = \alpha_j - \sum^{m}_{i=1} c^i*v_{i,j}\]
	\item 公开证明对 $(c, \boldsymbol{r})$。
\end{enumerate}


\subsubsection*{验证 (Verification)}

\begin{enumerate}
	\item 计算挑战：\(c' = \mathcal{H}(...,\textbf{r} \bullet \textbf{H} + \sum^{m}_{i=1} (c')^i*C_i)\)。
	\item 若 $c = c'$，则证明者必然知道 $\boldsymbol{V}$（概率可忽略不计的情况除外）。
\end{enumerate}

\subsubsection*{原理 (Why it works)}

\begin{align*}
\textbf{r} \bullet \textbf{H} &= C_{\alpha} - \sum^{m}_{i=1} (c')^i*C_i \\
\sum^{N}_{j=1} [r_j*H_j] &= C_{\alpha} - \sum^{m}_{i=1} (c')^i*C_i \\
\sum^{N}_{j=1} [(\alpha_j - \sum^{m}_{i=1} c^i*v_{i,j})*H_j] &= C_{\alpha} - \sum^{m}_{i=1} (c')^i*C_i \\
(\sum^{N}_{j=1} \alpha_j*H_j) - \sum^{N}_{j=1} [(\sum^{m}_{i=1} c^i*v_{i,j})*H_j] &= C_{\alpha} - \sum^{m}_{i=1} (c')^i*C_i \\
C_{\alpha} - \sum^{m}_{i=1} c^i*[\sum^{N}_{j=1} v_{i,j}*H_j] &= C_{\alpha} - \sum^{m}_{i=1} (c')^i*C_i \\
C_{\alpha} - \sum^{m}_{i=1} c^i*C_i &= C_{\alpha} - \sum^{m}_{i=1} (c')^i*C_i \\
\end{align*}

验证者可以确信证明者知道 $\boldsymbol{V}$ 的所有元素（概率可忽略不计的情况除外），这得益于本证明中多项特性的综合效果（结合了基本 Schnorr 签名的逻辑，见第 \ref{sec:schnorr-fiat-shamir} 节的说明）。

\subsubsection*{各特性的作用 (Justification of Features)}

\begin{enumerate}
    \item 如果我们不使用 $\boldsymbol{H}$ 或 $c$ 的幂次，响应改为 $r = \sum \boldsymbol{\alpha} + c* \sum_{i=1}^{m} (\sum_{j=1}^{N} v_{i,j})$，那么证明者最多只证明了他们知道所有向量元素之和。换言之，就只是 $K = \sum_{i=1}^{m} C_i$ 的离散对数，如同普通 Schnorr 签名。

    \item 加入 $\boldsymbol{H}$ 使 $r_j = \alpha_j +  c* \sum_{i=1}^{m} v_{i,j}$，这确保了证明者至少知道每个向量索引 $j$ 处的和。虽然可能存在许多其他 $\boldsymbol{r'}$ 向量使得 $R = \boldsymbol{r'} \bullet \boldsymbol{H} = r'_1*\lambda_1*G + r'_2*\lambda_2*G + ...$，但在不知道 $\boldsymbol{\lambda}$（我们假设其未知）的情况下，仅凭所有元素之和（步骤 1）很难找到其中任何一个。\footnote{即使两个或更多向量索引的和相同（$\sum_{i=1}^{m} v_{i,a} = \sum_{i=1}^{m} v_{i,b} = k$），$k$ 也不会被泄露，因为在方程对 $r_a = \alpha_a + c*k$ 和 $r_b = \alpha_b + c*k$ 中有三个未知数（$\alpha_a$、$\alpha_b$ 和 $k$），每增加一个方程 $r_q = \alpha_q + c*k$ 就多一个未知数（$\alpha_q$）。}

    \item 引入 $c$ 的幂次使每个响应元素 $r_j$ 成为 $c$ 的多项式函数，$r_j(c) = \alpha_j + v_{1,j}*c + v_{2,j}*c^2 + ... + v_{m,j}*c^m$。给定 $c$，存在许多 $(\alpha'_j, v'_{1,j}, ..., v'_{m,j})$ 的组合可以产生 $r_j$，但猜中的概率可以忽略不计，因为 $r_j$ 本身是未知的（可以是 1 到 $l$ 之间的任何值 [$l$ 是椭圆曲线子群的阶]）。即使每个向量索引处的和已知（步骤 2），随着 $\boldsymbol{V}$ 大小的增长，猜中有用的因式分解变得越来越困难（通常唯一可行的就是 $\boldsymbol{V}$ 的实际元素）。此外，$c$ 是事先未知的，因此证明者不能报告一个随机的 $r_j$（回忆第 \ref{sec:schnorr-fiat-shamir} 节）。这意味着证明者必须知道 $\boldsymbol{V}$ 的所有元素。\footnote{部分了解 $\boldsymbol{V}$ 的人可以提高伪造证明的概率（在大多数情况下仅提高可忽略不计的幅度）。一个极端的例子是，如果他知道了 $\boldsymbol{V}$ 中除一个元素以外的所有元素，他知道所有承诺 $C_j$ 及其盲化因子，并且知道缺失元素在范围（$q$ 到 $p$）内（其中 $p - q < l$），那么他猜中该元素（并通过计算 $C'_j$ 来验证）的概率高于通过猜测和检验来求解离散对数问题。不过在这种情况下他最终还是知道了 $\boldsymbol{V}$ 的所有元素。这一逻辑也适用于其他部分知识场景。}
\end{enumerate}


%----------
\section{内积证明 (Inner Product Proof)}

假设有两个向量 $\boldsymbol{v}$ 和 $\boldsymbol{z}$，各有 $N$ 个元素。它们的内积为 $s = \boldsymbol{v} \bullet \boldsymbol{z}$，我们有它们的承诺 $[C_s = x_s G + s H_1, C_v = (x_v, \boldsymbol{v}) \bullet (G,\boldsymbol{H}), C_z = (x_z, \boldsymbol{z}) \bullet (G,\boldsymbol{H})]$，\footnote{此处我们的记号 $(x_v, \boldsymbol{v}) \bullet (G,\boldsymbol{H})$ 意味着盲化因子 $x_v$ 被放在 $\boldsymbol{v}$ 前面进行内积运算，生成元 $G$ 同理。}我们希望证明内积方程成立（并且我们知道所有元素），同时不泄露关于 $(s, \boldsymbol{v}, \boldsymbol{z})$ 的任何信息（除了 $N$）。换言之，证明 $C_s$ 中的值是 $C_v$ 和 $C_z$ 中大小为 $N$ 的向量的内积。注意现在我们引入了盲化因子 $x_s$、$x_v$ 和 $x_z$，它们严格来说不属于被考虑的向量。这一概念在后续章节中很有用，因为用随机掩码隐藏原始值非常重要。

基本上，我们对 $\boldsymbol{v}$ 和 $\boldsymbol{z}$ 进行两次向量知识证明（第 \ref{sec:vectorzkproof} 节），并额外做一次内积证明。


\subsubsection*{非交互式证明 (Non-interactive proof)}

\begin{enumerate}
	\item 生成两个向量（大小 N）和四个整数 \((\alpha_{v,0}, \boldsymbol{\alpha_v}, \alpha_{z,0}, \boldsymbol{\alpha_z}, \alpha_{s,0}, \alpha_{s,1}) \in_R \mathbb{Z}_l\)，并计算
	\begin{align*}
	C_{\alpha}^{v} &= (\alpha_{v,0},\boldsymbol{\alpha_v}) \bullet (G,\textbf{H}) \\
	C_{\alpha}^{z} &= (\alpha_{z,0}, \boldsymbol{\alpha_z}) \bullet (G,\textbf{H}) \\
	C_{\alpha}^{s,0} &= \alpha_{s,0}*G + (\boldsymbol{\alpha_v} \bullet \boldsymbol{\alpha_z})*H_1 \\
	C_{\alpha}^{s,1} &= \alpha_{s,1}*G + (\boldsymbol{v} \bullet \boldsymbol{\alpha_z} + \boldsymbol{z} \bullet \boldsymbol{\alpha_v})*H_1
	\end{align*}
	\item 计算{\em 挑战} $c = \mathcal{H}(...,C_{\alpha}^{v},C_{\alpha}^{z},C_{\alpha}^{s,0},C_{\alpha}^{s,1})$。
	\item 定义{\em 响应} $\boldsymbol{r}$，包含向量（大小 N）$\boldsymbol{r_v}, \boldsymbol{r_z}$，其中 $r_{v,0}, r_{z,0}$ 对应向量的盲化因子，$r_s$ 对应内积的盲化因子
	\begin{align*}
	r_{v,0} &= \alpha_{v,0} - c*x_v \\
	r_{z,0} &= \alpha_{z,0} - c*x_z \\
	r_{v,j} &= \alpha_{v,j} - c*v_j \\
	r_{z,j} &= \alpha_{z,j} - c*z_j \\
	r_s &= \alpha_{s,0} + c*\alpha_{s,1} + c^2*x_s 
	\end{align*}
	\item 公开证明 $(c, \boldsymbol{r}, C^{s,0}_{\alpha}, C^{s,1}_{\alpha})$。
\end{enumerate}

\subsubsection*{验证 (Verification)}

\begin{enumerate}
	\item 计算挑战：
	\[c' = \mathcal{H}(...,[(r_{v,0},\boldsymbol{r_v}) \bullet (G,\textbf{H}) + c*C_v],[(r_{z,0},\boldsymbol{r_z}) \bullet (G,\textbf{H}) + c*C_z],[C^{s,0}_{\alpha}],[C^{s,1}_{\alpha}])\]
	\item 计算\footnote{由于 $C^{s,0}_{\alpha}$ 和 $C^{s,1}_{\alpha}$ 是关联的，据我们所知无法将 $R_s$ 移入挑战的计算中。}
	\begin{align*}
	%R_v &= (r_{v,0},\boldsymbol{r_v}) \bullet \textbf{H} \\
	%R'_v &= C_{\alpha}^{v} + c'*C_v \\
	%R_z &= (r_{z,0},\boldsymbol{r_z}) \bullet \textbf{H} \\
	%R'_z &= C_{\alpha}^{z} + c'*C_z \\
	R_s &= r_s*G + (\boldsymbol{r_v} \bullet \boldsymbol{r_z})*H_1 \\
	R'_s &= C_{\alpha}^{s,0} + c'*C_{\alpha}^{s,1} + c'^2*C_s
	\end{align*}
	\item 若 $c = c'$ 且 $R_s = R'_s$，则证明者必然知道 $\boldsymbol{v}, \boldsymbol{z}$ 和 $s$，且内积 $s = \boldsymbol{v} \bullet \boldsymbol{z}$ 成立（概率可忽略不计的情况除外）。
\end{enumerate}

\subsubsection*{原理（内积部分）(Why it works (inner product component))}

\begin{align*}
r_s*G + (\boldsymbol{r_v} \bullet \boldsymbol{r_z})*H_1 &= C_{\alpha}^{s,0} + c*C_{\alpha}^{s,1} + c^2*C_s \\
(\alpha_{s,0} + c \alpha_{s,1} + c^2 x_s)*G + (\boldsymbol{\alpha_v} \bullet \boldsymbol{\alpha_z} + c[\boldsymbol{v} \bullet \boldsymbol{\alpha_z} + \boldsymbol{z} \bullet \boldsymbol{\alpha_v}] + c^2 [\boldsymbol{v} \bullet \boldsymbol{z}])*H_1 &= C_{\alpha}^{s,0} + c*C_{\alpha}^{s,1} + c^2*C_s \\
C_{\alpha}^{s,0} + c*C_{\alpha}^{s,1} + c^2 (x_s*G + [\boldsymbol{v} \bullet \boldsymbol{z}]*H_1) &= C_{\alpha}^{s,0} + c*C_{\alpha}^{s,1} + c^2*C_s \\
c^2 (x_s*G + [\boldsymbol{v} \bullet \boldsymbol{z}]*H_1) &= c^2(x_s*G + s*H_1) \\
\boldsymbol{v} \bullet \boldsymbol{z} &= s
\end{align*}

响应 $\boldsymbol{r_v}$ 和 $\boldsymbol{r_z}$ 证明了对向量 $\boldsymbol{v}$ 和 $\boldsymbol{z}$ 的知识，因此在 $R_s \stackrel{?}{=} R'_s$ 中使用它们即可证明内积成立。读者可以自行探索代数逻辑来确认这一点。


\section{压缩向量知识证明 (Condensed Vector Knowledge Proof)}
\label{sec:condensedvectorproof}

我们最初的向量知识证明为 $(c, \boldsymbol{r})$，其证明大小与向量维度呈线性关系。若 $N = 500$（且 $m = 1$），证明将占用 $(1 + 500)*32 = 16032$ 字节。我们可以使用 Bootle 等人提出的方法 \cite{bootle-efficient-zkcircuits} 来压缩证明大小。它与向量维度呈对数关系（证明大小 $\approx 6*log_2(N)$，因此对于 N = 500，证明大小约为 54*32 = 1728 字节）。


\subsection*{非交互式证明 (Non-interactive proof)}

\begin{enumerate}
    \item 以最小化证明大小为目标（优化验证时间更为复杂），将 $N$ 分解为 $q$ 个素数因子，按从大到小排列。若 $N = 500$，则 $\boldsymbol{f} = \langle 500, 5, 5, 5, 2, 2 \rangle$，索引从 $0$ 到 $q = 5$（第 0 项为原始 N）。
    \item 对 $i = 1$ 到 $q$（索引 0 对应原始向量），
    \begin{enumerate}
        \item 将 $\boldsymbol{v}^{i-1}$ 和 $\boldsymbol{G}^{i-1}$ 按大小 $f[i]$ 分块为更小的向量
        \[ \langle v^{i-1}[1], ..., v^{i-1}[f[i-1]] \rangle \xrightarrow{} \langle \langle v^{i-1}[1],...,v^{i-1}[f[i] \rangle,\langle \rangle,...,\langle v^{i-1}[f[i-1] - f[i] ],..., v^{i-1}[f[i-1]\rangle \rangle \]
    \end{enumerate}
    \[\begin{pmatrix}
        1 & 2 & 3 \\
        4 & 5 & 6
    \end{pmatrix}\]
\end{enumerate}


\subsection*{验证 (Verification)}
